\section{Pourquoi notre nouvel environnement?}

Tout d'abord, m\^eme si l'histoire de l'\'etude des AC est presque aussi
ancienne que l'histoire des ordinateurs~\cite[]{Metropolis93b}, les
architectures cellulaires construites, bien que non encore
industrialis\'ees apparaissent comme les plus prometteuses\footnote{%
  Il est difficile de savoir si ce fait est \`a mettre au cr\'edit des
  architectures cellulaires, nous nous bornerons \`a proposer l'id\'ee que
  la perception du potentiel d'usage d'une architecture est
  inversement proportionnelle \`a sa r\'ealisation industrielle.}, %
et le probl\`eme des m\'ethodes de constructions d'architectures et
d'algorithmes cellulaires restent largement ouvert. %
De m\^eme, \`a propos du probl\`eme de
l'auto-reproduction~\cite[]{vonNeumann66}, la dualit\'e opposant
analytique \`a synth\'etique accompagne la naissance des AC\footnote{%
  Voir \cite{Reggia93} pour une approche permettant d'accorder plus
  d'autonomie aux ph\'enom\`enes d'auto-reproduction. %
  Nous consid\'erons, dans ce contexte, l'autonomie comme la mesure
  d'une position sur l'axe qui va du construit \`a l'\'emergeant.}. %
Cette dualit\'e, ou plut\^ot la capacit\'e offerte par les AC de basculer
d'un univers \'epist\'emologique \`a un autre tout en restant au sein d'un
mod\`ele unique, doit \^etre incarn\'ee dans l'environnement. %
Nous souhaitons donc un environnement permettant une approche mixte. %

Par ailleurs, aucun des outils que nous avons cit\'es ne r\'epond
totalement \`a un ensemble de contraintes plus pratiques que nous avons
fix\'ees comme cadre de d\'eveloppement. %
\begin{itemize}
\item la nature publique\footnote{%
    Par \aconcept{publique} nous entendons accessible sous forme
    source et utilisable sans licence commerciale.} %
  des outils
\item la rapidit\'e d'ex\'ecution
\item la richesse de visualisation
\item la souplesse de sp\'ecification
\item un environnement d'exploitation Unix et X11
\item la modularit\'e des constituants en vue de la construction et du
  s\'equencement d'exp\'eriences cellulaires. %
\end{itemize}

\asoft{CA~LAB} n'est pas un outil public et n'est pas disponible dans
l'environnement Unix. %
\asoft{Cellsim} de par sa d\'ependance \`a~\asoft{SunView} n'est
utilisable que dans un environnement Sun\footnote{%
  Felicity George a r\'ealis\'e en 1992 un portage de~\asoft{cellsim} pour
  l'environnement~\asoft{XView}.}, %
\asoft{scamper} pr\'esente une application int\'egrant un formalisme de
sp\'ecification des r\`egles bas\'ees par pattern, des compteurs de
populations et une interface visuelle bas\'ee sur X11, mais il ne permet ni
ex\'ecution rapide, ni s\'equencement. %
\asoft{Cellang} propose un langage de description d'AC, mais son
compilateur ne produit pas un code permettant de tirer le meilleur
parti des ressources disponibles. %

\section{Des AC pour l'optimisation d'usage de ressources}

Concernant le parall\'elisme massif, les tendances de la pr\'ec\'edente
d\'ecennie en architecture des machines ne confirment pas les
pr\'edictions avanc\'ees \`a son propos\footnote{%
  Machines SIMD massives reposant sur des techniques de fabrication de
  type wafer scale int\'egration, associ\'ees \`a des algorithmes de
  reconfigurations de r\'eseaux d'interconnexion permettant la
  r\'esistance aux d\'efauts; pyramides de telles architectures pour la
  vision et plus g\'en\'eralement l'IA, ou architectures diverses pour la
  physique, la biologie mol\'eculaire, le pattern matching haut d\'ebit,
  la r\'ealisation de film holographique temps r\'eel.}. %
N\'eanmoins, hors consid\'erations architecturales, les ressources
accessibles continuent \`a cro\^itre \`a un rythme soutenu\footnote{%
  Au point qu'on pourrait mettre en doute l'application du
  principe des gaz parfaits \`a la consommation de ressources. %
  Contre cette tendance, nous proposons de consid\'erer qu'il n'est plus
  temps d'apprendre comment se contenter de ressources restreintes,
  mais comment tirer profit de ressources disponibles.}. %
Hormis le calcul intensif classique, les gains de ressources semblent
s'orienter vers le parall\'elisme MIMD massif pour la gestion et le
support multim\'edia pour le poste de travail. %
Nous affirmons que les AC constituent un mod\`ele efficace et universel
pour l'utilisation de ces ressources, y compris par la prise en compte
du parall\'elisme multi-calculateur\footnote{%
  C'est \`a dire l'utilisation d'un ensemble de machines
  autonomes~(stations de travail ou autres, localement parall\`eles ou
  non) reli\'ees par un r\'eseau d'interconnexion~(de type bus ou point
  \`a point) dot\'ees d'une couche logicielle syst\`eme ad hoc, comme une
  ressource parall\`ele en soi.} %
associ\'e \`a l'interconnexion par \astrangeterm{WAN}{%
  Wide Area Network, il s'agit ici de remarquer que les d\'ebits
  disponibles ou annonc\'es pour les r\'eseaux de t\'el\'ecommunication publics
  permettent d'envisager la construction virtuelle de
  machines cellulaires reparties sur une large zone g\'eographique.} %
de ceux-ci. %
La r\'egularit\'e structurelle des donn\'ees et des traitements des AC
associ\'es \`a l'absence de probl\`eme algorithmique de synchronisation
des \'el\'ements du calcul r\'eparti, permet un calcul et une affectation de
ressource simple. %
Une approche bottom-up de la sp\'ecification du comportement des
\'el\'ements cellulaires permet, de plus, de s'affranchir des probl\`emes
d'architecture CPU, une fonction consid\'er\'ee sous l'angle de sa table
de transition est ind\'ependante des architectures de
traitement\footnote{%
  A une permutation de l'ordre des bits sur le mot pr\`es, et toujours en
  consid\'erant que l'unit\'e de traitement r\'eellement disponible reste
  fondamentalement une architecture de von Neumann dot\'ee d'une m\'emoire
  adressable importante.}. %

\mysetref{repartition-du-calcul}
A propos de la r\'epartition d'un calcul cellulaire sur r\'eseau de
stations de travail\footnote{%
  Algorithme distribu\'e avec recouvrement: calcul des bords, envoi des
  bords, calcul du centre, r\'eception des bords.}, %
et concernant le probl\`eme de l'overhead associ\'e au passage de messages
impl\'ementant la synchronisation des bords lorsque le calcul courant
r\'eclame un acc\`es \`a une cellule voisine, ind\'ependamment d'une mesure
particuli\`ere des d\'ebits respectifs du couple
\bialt{cpu}{m\'emoire}~($D_{m}$) et du canal de communication($D_{c}$),
le rapport des tailles \bialt{locales}{distantes} $N^{d}/N^{d-1} = N$
garantit l'existence d'une valeur $k$ de $N$ telle que l'usage de
plusieurs n{\oe}uds de calculs sera plus efficace que l'usage d'un
seul d'entre eux:~$\frac{D_{m}}{D_{c}} k > 1$. %
%% \[\frac{D_{m}k^{d}}{D_{c}k^{d-1}} > 1\]

A d\'efaut de machines cellulaires industrielles, nous
esp\'erons pouvoir utiliser les DSP et les acc\'el\'erateurs
\aquote{bitblt}~\cite[]{Pike84} aujourd'hui largement disponibles
comme acc\'el\'erateurs de nos moteurs cellulaires. %

\section[AC et morphog\'en\`ese]{Des AC comme outils de r\'eflexion morphog\'en\'etique}

L'essence de la forme pose un probl\`eme, ne serait ce que
par l'existence de la g\'en\'ericit\'e du moule. %
Tout \'etat du monde n'est pas une forme, toutefois une forme est par
d\'efinition une configuration d'espace cellulaire prise parmi un
ensemble fini de configurations. %
Cette d\'efinition r\'eclame un crit\`ere de
\bialt{discrimination}{r\'eduction} de l'ensemble des \'etats. %
Une premi\`ere possibilit\'e r\'eside dans l'utilisation d'une fonction
arbitraire sur un \'etat d\'eterminant la valeur d'un attribut formel:
esth\'etique, g\'eom\'etrique, sta\-tistique, arithm\'etique. %
Les formes sont figuratives, abstraites, s\'eduisantes, rondes, carr\'ees,
sym\'etriques, rares, uniques\footnote{%
  Chaque configuration d'espace est \'evidemment unique, on pourrait
  n\'eanmoins mesurer, intuitivement ou formellement, l'unicit\'e \`a partir
  du produit de l'ensemble des \'etats par l'ensemble des
  fonctions de qualification formelle disponibles. %
  Les fonctions de qualification sont des fonctions de r\'eduction de
  l'information~(le nombre de bits rendu est inf\'erieur au nombre de
  bits fourni), le cardinal de l'image d'un ensemble de formes est
  donc inf\'erieur \`a celui de l'ensemble de formes, la distribution des
  redondances de valeurs induites par la diff\'erence de cardinalit\'e
  d\'eterminant les classes d'\'equivalence de cet attribut. %
  
  Une relation d'ordre pour l'unicit\'e est alors d\'efinie par une
  fonction minimisant le cardinal de l'union des classes d'\'equivalence
  des images des formes.}, %
r\'eguli\`eres, fractales, denses, color\'ees, leur num\'ero d'ordre enfin,
peut \^etre argument d'une fonction choisie au hasard. %
Cette m\'ethode pose le probl\`eme de l'\'enum\'eration des fonctions de
r\'eduction\footnote{%
  Ainsi que celui du choix \`a priori des \'etats consid\'er\'es. %
  Les fonctions de qualification sont implicitement destin\'ees \`a
  introduire une r\'eduction du nombre des formes par l'introduction de
  classes d'\'equivalence sur l'ensemble des \'etats.}. %
Une deuxi\`eme m\'ethode consiste \`a opposer la g\'en\'eticit\'e \`a la
g\'en\'ericit\'e, c'est \`a dire l'histoire de la forme au moule de la forme. %
Dans ce contexte la r\'eduction du nombre d'espaces appel\'es \`a la forme
est d\'etermin\'ee par l'existence d'une trajectoire dans l'espace des
\'etats. %
Cette trajectoire \'etant elle m\^eme d\'etermin\'ee par l'application it\'er\'ee
d'une fonction de l'ensemble des \'etats vers l'ensemble des \'etats. %
Les trajectoires ne sont pas moins nombreuses que les \'etats, mais si
l'on ajoute une contrainte de localit\'e aux fonctions g\'en\'eratrices de
trajectoires, une tentative d'\'enum\'eration devient accessible. %
Cette contrainte de localit\'e constitue le fondement d'une approche
cellulaire de la forme. %
Cette contrainte introduit un lien entre la vitesse et la forme. %
Plus g\'en\'eralement, la notion de contrainte s'applique ais\'ement au
monde des formes dans une perspective esth\'etique, le libre choix des
contraintes de production de l'objet d\'eterminant la qualit\'e artistique
de l'{\oe}uvre, dont la forme peut \^etre con\c{c}ue comme une trace
d'ex\'ecution\footnote{%
  La perception de l'{\oe}uvre n'est jamais exempte de la perception
  de l'histoire de sa fabrication~\cite[]{Gombrich71}.}. %
De m\^eme la \aquote{forme} d'un espace cellulaire
devient la trace d'un programme\footnote{%
  De l'application d'une loi non ergodique sur un espace initial
  n\'ecessairement al\'eatoire.}. %

\section[AC, espace, temps et calcul]%
{Des AC comme outils de r\'eflexion sur l'espace, le temps et le calcul}

Nous avons d\'ej\`a mentionn\'e l'existence d'une dualit\'e
analytique-synth\'etique. %
D'autres projections de cette dualit\'e sont possibles, sur les axes
dynamique statique, heuristique algorithmique, parall\`ele s\'equentiel ou
temporel spatial. %
La trajectoire d'un AC sur un espace de configuration est un objet
dynamique, la structure d'un espace de r\`egles de transition d'AC pour
un voisinage donn\'e est un objet statique. %
A la fronti\`ere se trouvent des objets donn\'es par l'application
it\'erative d'une fonction simple sur un espace cellulaire initial
globalement param\'etr\'e~(l'automate auto-reproducteur de von Neuman, ou
une machine de Turing impl\'ement\'ee par un espace de configuration
d'un AC tel \arule{Life} de Conway). %

La quantit\'e de parall\'elisme d'une ex\'ecution peut se mesurer au nombre
d'instructions diff\'erentes rencontr\'ees lors de cette ex\'ecution. %
Il existe une \'equivalence entre un AC d\'efini par un r\'eseau
d'interconnexion et une fonction de transition, et un algorithme
cellulaire d\'efinit par une s\'equence d'instructions
cellulaires. %
Si un AC est consid\'er\'e comme une instruction cellulaire \'el\'ementaire,
celle-ci poss\`ede un co\^ut, qui s'exprime comme la taille de la table
de v\'erit\'e de cette instruction cellulaire. %
Si un AC est consid\'er\'e comme une s\'equence d'instructions cellulaires,
celle-ci poss\`ede un co\^ut, qui s'exprime comme la somme des temps de
r\'ealisation de chaque instruction, et in fine la somme de la taille
des tables de v\'erit\'e de chaque instruction \'el\'ementaire. %
L'impl\'ementation optimale d'un algorithme cellulaire serait obtenu
lorsque, pour une occupation maximale de la m\'emoire disponible pour les
tables de v\'erit\'es~(les fonctions de transition) des instructions
cellulaires, le nombre de primitives diff\'erentes composant l'algorithme
est minimum. %
Le ratio du nombre de primitives employ\'ees sur le nombre de primitives
existantes d\'etermine l'importance des effets relevant de la
dynamique~(l'heuristique) de ceux relevant de l'algorithme proprement
dit dans la fonction associ\'ee au programme cellulaire.

Nous disons que la fonction de transition d'un AC est localement
univer\-selle\footnote{%
  Nous dirons parfois qu'un AC est globalement universel, pour
  d\'esigner la propri\'et\'e usuelle d'universalit\'e par opposition \`a la
  propri\'et\'e d'universalit\'e locale des PE de machines SIMD\@. %
  L'universalit\'e locale implique \'evidement l'universalit\'e globale. %
  \arule{Life} constitue dans le contexte de ce travail l'arch\'etype
  d'un AC dot\'e de l'universalit\'e globale, et le processeur GAPP
  l'arch\'etype d'un AC dot\'e de l'universalit\'e locale.} %
si une seule cellule consid\'er\'ee comme PE se comporte comme une machine
universelle sous r\'eserve d'une extension virtuelle du nombre d'\'etats
cellulaires. %
C'est a priori le cas de toutes les architectures SIMD,
m\^eme si le faible nombre d'\'etats accessibles\footnote{%
  Ce faible nombre d'\'etats doit n\'eanmoins \^etre exprim\'e en logarithme,
  c'est \`a dire en nombre de bits adresse.} %
destine \'evidemment ces machines \`a un usage cellulaire. %

D'un cot\'e le \aquote{calcul s\'equentiel} pur d\'eroule une s\'equence
d'instructions d\'efinie en extension sur un espace initial \'egalement
d\'efini en extension, de l'autre le \aquote{calcul cellulaire pur}
it\`ere une loi unique d\'efinie en extension sur un espace initial dont
seules des propri\'et\'es statistiques\footnote{%
  La taille d'un espace cellulaire incarn\'e, n\'ecessairement fini, doit
  n\'eanmoins \^etre consid\'er\'ee comme tendant vers l'infini; la
  sp\'ecification d'une configuration cellulaire particuli\`ere en
  extension devient alors une barri\`ere \`a la pleine utilisation des
  ressources disponibles.} %
sont d\'efinies. %
La distinction entre une s\'equence d'instructions et une \emph{loi}
apparaissant comme une instruction unique, ne tient ni \`a la taille de
la loi ni \`a une propri\'et\'e telle l'universalit\'e de l'architecture
sous-tendant la s\'equence d'instructions, ou l'universalit\'e
locale ou globale de la loi\footnote{%
  En effet toute s\'equence d'instruction, tout algorithme peuvent
  \'evidemment \^etre transform\'es en une instruction
  unique par le calcul d'une lookup-table.} %
mais \`a l'opposition entre l'unicit\'e de l'application d'une s\'equence
d'instructions, d'un algorithme et la n\'ecessaire it\'eration sans
fin\footnote{%
  Comme pour l'espace, ressource n\'ecessairement limit\'ee devant
  n\'eanmoins \^etre consid\'er\'ee comme tendant vers l'infini, le temps
  cellulaire, le nombre de cycles d'it\'erations de l'application de la
  loi sur l'espace doit \^etre consid\'er\'e comme tendant vers l'infini.} %
de la loi sur l'espace cellulaire. %
A la fronti\`ere\footnote{%
  La structure de l'espace englobant le s\'equentiel et le cellulaire
  appara\^it ici tr\`es clairement.} %
le calcul cellulaire it\'erant une loi unique sur un espace dont la
configuration ini\-tiale est une transformation~(un mapping) d'une s\'erie
d'instructions d'un calcul s\'equentiel, c'est \`a dire le calcul
cellulaire exploitant l'universalit\'e globale d'une loi; ou encore le
calcul cellulaire SIMD d\'eroulant une s\'equence d'instructions d\'efinie
en extension, c'est \`a dire le calcul cellulaire utilisant
l'universalit\'e locale d'une loi, et sans doute, mais sans qu'il soit
possible d'en \'etablir simplement la part, l'universalit\'e globale de
cette loi. %

Si l'on consid\`ere que l'encodage d'un algorithme sur une machine
s\'equentielle pose un probl\`eme aigu, et encore davantage
l'encodage du probl\`eme \`a traiter, il en va bien s\^ur de m\^eme pour les
machines SIMD. %
On continue toutefois dans le contexte SIMD \`a r\'esoudre des probl\`emes
algorithmiques, \`a utiliser l'universalit\'e pour la construction de
machines algorithmes. %
Lorsque l'on passe \`a une machine cellulaire pure, la nature
intrins\`equement statistique du processus de calcul\footnote{%
  Nous avons jusque l\`a \'evoqu\'e la seule incertitude statistique
  associ\'ee \`a la sp\'ecification de la configuration initiale, la loi
  \'etant consid\'er\'ee comme purement d\'eterministe, on voit bien cependant
  comment un ind\'eterminisme sur l'espace est de fait indiff\'erentiable
  d'un ind\'eterminisme sur la loi, laquelle n'est jamais que de
  l'espace local g\'en\'erateur du temps.} %
change radicalement le point de vue. %
Le probl\`eme \`a traiter ne pouvant d\`es lors gu\`ere s'exprimer que par
analogie \`a une morphog\'en\`ese
cellulaire\footnote{%
  Expression \`a prendre dans un sens non restrictif \`a la biologie bien
  s\^ur, \`a l'existence pr\`es d'un paysage \'epig\'en\'etique.}. %
Il sera toujours possible de plonger une approche SIMD, voire
purement s\'equentielle et distribu\'ee dans un contexte cellulaire pur,
mais outre que cela pose des probl\`emes qui nous apparaissent assez
ardus de faisabilit\'e pratique\footnote{%
  Cela pose aussi des probl\`emes plus fondamentaux, dont le moindre
  n'est pas d'\'eclaircir la nature ontologique de l'incertitude
  statistique qui semble bien devoir \^etre li\'ee \`a une d\'efinition d'un
  calcul cellulaire r\'eellement massif, m\^eme si la dimension
  technologique de cette incertitude existe \'egalement.}, %
cela ne tend pas \`a r\'epondre \`a la
question de la sp\'ecificit\'e du calcul cellulaire pur. %
En effet, si celui-ci ne semble pas, dans son approche la plus stricte,
isomorphe \`a une machine universelle arbitraire, ce pourrait \^etre d\^u \`a
la nature particuli\`ere des morphog\'en\`eses cellulaires, lesquelles ne
sont peut-\^etre pas r\'eductibles \`a des algorithmes. %
Le calcul s\'equentiel peut \^etre consid\'er\'e comme une version
adimensionnelle~(de dimension z\'ero) d'un calcul cellulaire. %
Le calcul cellulaire algorithmique\footnote{%
  C'est \`a dire la r\'ealisation d'un algorithme par une une
  configuration cellulaire initiale associ\'ee \`a une loi globalement
  universelle.} %
correspondant au choix d'un point unique dans l'espace des \'etats
cellulaires, permettant par le plongement d'une machine de Turing dans
l'espace cellulaire, de r\'ealiser une r\'eduction
dimensionnelle de l'espace cellulaire. %
Le calcul cellulaire SIMD\footnote{%
  C'est \`a dire la r\'ealisation d'un algorithme sur une configuration
  cellulaire par une m\'ethode s\'equentielle locale.} %
correspondant au choix d'un mapping probl\`eme vers configuration cellulaire,
associ\'e ou non \`a l'utilisation des propri\'et\'es dynamiques d'une loi ou
plusieurs lois pour contr\^oler une trajectoire dans l'espace des
configurations cellulaires. %
Le calcul cellulaire pur, par la prise en compte de l'incertitude
statistique, cherche \`a dresser un paysage g\'en\'eral du calcul,
lequel devient une propri\'et\'e intrins\`eque, voire \'emergeante d'un mod\`ele
d'univers compos\'e d'atomes d'espace dot\'e d'une loi universelle. %
Une des propri\'et\'es particuli\`erement int\'eressante, de notre point de
vue, du calcul cellulaire, r\'eside dans les \'equivalences espace-temps
rendues possibles par son architecture. %
En effet la mesure des quantit\'es \'el\'ementaires de cette architecture,
ainsi que les transformations mutuelles possibles entre ces quantit\'es
sont d'impl\'ementation ais\'ee. %
Ainsi, en calcul cellulaire pur, le nombre d'\'etats d'une cellule et le
nombre de voisins de cette cellule d\'etermine une borne sup\'erieure \`a la
complexit\'e\footnote{%
  Nous n'avons pas de d\'efinition pr\'ecise de cette complexit\'e,
  il s'agit a priori d'une mesure statistique de la richesse de
  comportement, rapport\'ee \`a une dur\'ee. %
  La capacit\'e et la rapidit\'e morphog\'en\'etique associ\'ees \`a une r\`egle.} %
du calcul effectu\'e\footnote{%
  L'analyse de la table de transition permet de pr\'eciser une valeur
  de complexit\'e, en particulier la distance \`a la borne sup\'erieure. %
  Une table constante, les tables identit\'e ou inverse sur un seul
  bit du mot de voisinage~(translation, et translation alternance de
  la configuration initiale) sont faciles \`a identifier. %
  Une g\'en\'eralisation possible de ces notions conduit \`a la d\'efinition
  de param\`etres de contr\^ole de l'espace des r\`egles tels que le
  param\`etre $\lambda$ de Langton ou le param\`etre d'ench\^assement.}. %
La taille de \arule{Life} est de~512~bits,
celle de l'automate auto-reproducteur de von Neumann de
$2^{5\lceil\log_{2}(29)\rceil}\lceil\log_{2}(29)\rceil =
20 \text{m\'egaoctets}$. %
Il serait int\'eressant de conna\^itre la taille minimale d'une
configuration initiale de \arule{Life} r\'ealisant une machine
universelle, il n'existe pas \`a notre connaissance de r\'ealisation
effective d'une telle configuration. %
La taille de la configuration r\'ealisant l'automate auto-reproducteur
de von Neumann a \'et\'e \'evalu\'ee \`a plusieurs centaines de milliers de
cellules. %
%% John kemeny 1955
Jaqueline Signorini a \'etudi\'e les conditions de r\'ealisation de cet
automate sur le GAPP~\cite[]{Signorini92}, mais une configuration
auto-reproductrice compl\`ete n'a encore jamais \'et\'e men\'ee \`a bien. %

L'impl\'ementation de \arule{Life} sur une machine s\'equentielle
$M_{\text{SISD}}$ utilise $a$ cycles de base pour chacune des $N$
cellules, et $b = 2$ bits\footnote{% 
  Un plan \avar{past} et un plan \avar{futur}, impl\'ementant le
  synchronisme de mise \`a jour, on peut envisager des optimisations
  r\'eduisant le nombre de bits n\'ecessaire au plan secondaire,
  lesquelles se ram\`enent \`a une compression de la configuration
  cellulaire, associ\'ee \`a des m\'ethodes d'acc\`es sp\'ecialis\'ees au codage de
  la compression, renfor\c{c}ant ainsi la tendance des machines
  s\'equentielles \`a l'\'echange de temps contre de l'espace.} %
par cellule pour le calcul d'une g\'en\'eration de \arule{Life}, plus
$c \times  i$ bits pour le codage des instructions, $i$ \'etant la taille
moyenne d'une instruction. %
$c$ et $a$ sont typiquement tous les deux de l'ordre de quelques
centaines d'\bialt{instructions}{cycle} \'el\'ementaire d'une machine
monoprocesseur id\'eale pour une version na\"ive~(non optimis\'ee) de
l'\'emulateur de la machine cellulaire. %
Une impl\'ementation par lookup-table avec collecte pipelin\'ee du
voisinage augmente $c$ et diminue $a$ d'un ordre de grandeur. %

\begin{equation}
  \frac{\text{Space}}{\text{Time}} = \frac{bN+ci+M_{\text{SISD}}}{aN}
  \label{eq-life-sisd}
\end{equation}

Avec:\[b = 2,~i = 32,~c \approx a \approx 10^2,~M_{\text{SISD}}
\approx 10^6\]

La valeur de $a$ et $c$ d\'ependent du nombre $k$ d'\'etats par cellule et
du rayon $r$ du voisinage cellulaire. %
La taille du mot de voisinage~(9~dans le cas de \arule{Life})
d\'etermine ces valeurs. %

L'impl\'ementation de~\arule{Life} sur une machine
SIMD~$M_{\text{SIMD}}$ utilise $a^{\prime}$~cycles de base pour le
calcul d'une g\'en\'eration, ind\'ependamment du nombre de cellules,
et~$b^{\prime}$~bits par cellule, les~$c^{\prime} \times  i^{\prime}$~bits
d'instructions \'etant traditionnellement consid\'er\'es comme n'appartenant
pas stricto sensu \`a la machine SIMD, mais \`a une indispensable machine
s\'equentielle, h\^ote de toute machine SIMD\@. %
Les $a^{\prime}$~cycles de calcul d\'ependent de la machine SIMD choisie,
une machine SIMD typique est plus difficile \`a proposer qu'une
machine SISD id\'eale, sur le GAPP, de l'ordre de la centaine pour une
impl\'ementation a priori de la fonction de transition\footnote{%
  L'algorithme arithm\'etique de \arule{Life}, comptage plus seuil.}, %
de l'ordre du millier pour une impl\'ementation a posteriori de celle-ci\footnote{%
  Une d\'erivation algorithmique du DFA de \arule{Life} calcul\'ee \`a partir
  de la table de transition, laquelle peut \^etre obtenue par observation
  d'une trajectoire sur l'espace des configurations, associ\'ee \`a une
  algorithmique SIMD g\'en\'erique interpr\'etant le DFA sur le mot de voisinage.}. %
De l'ordre d'une dizaine de bits pour les $b^{\prime}$ bits d'espace. %

\begin{equation}
  \frac{\text{Space}}{\text{Time}} = \frac{(b^{\prime}+M_{\text{SIMD}})N}{a^{\prime}}
  \label{eq-life-simd}
\end{equation}

Avec:\[b \approx 10,a^{\prime} \approx 10^2,~\text{SIMD} \approx 10^2\]

L'impl\'ementation de \arule{Life} sur une machine cellulaire pure $M_{\text{CA}}$
utilise $a^{\prime\prime} = 1$ cycle et  $b^{\prime\prime} = 1$ bit par
cellule. %
Par contre le co\^ut li\'e \`a l'architecture s'exprime simplement comme la
taille de la table de transition. %

\begin{equation}
  \frac{\text{Space}}{\text{Time}} =
  \frac{(b^{\prime\prime}+M_{\text{CA}})N}{a^{\prime\prime}}
  \label{eq-life-cell}
\end{equation}

Avec:\[a^{\prime\prime} = 1,~b^{\prime\prime} = 1,~\text{CA} = 512\]

Les valeurs d'espace propos\'ees pour le terme $M_{\text{SISD}}$ et le
facteur $M_{\text{SIMD}}$ dans les \'equations~\ref{eq-life-sisd}
et~\ref{eq-life-simd} repr\'esentent un nombre de transistors CMOS id\'eal
pour ces architectures. %
Le facteur $M_{\text{CA}}$ de l'\'equation~\ref{eq-life-cell} n'est pas
un nombre de transistors approxim\'e, mais un nombre exact de bits. %
Cette diff\'erence de type entre les objets compt\'es se trouve
consid\'erablement pond\'er\'ee par les croissances respectives des tailles
d'espaces n\'ecessaires \`a la r\'ealisation d'une fonction de transition
lorsque la fonction de transition devient un param\`etre des \'equations
pr\'ec\'edentes. %
Sur une machine SISD l'augmentation de la taille du mot de voisinage
induit une croissance lin\'eaire de la taille de l'espace et du temps
utilis\'e pour calculer une trajectoire cellulaire. %
L'augmentation de la taille d'espace ne d\'epend que de la taille du mot
de voisinage, pas de la complexit\'e de la fonction de
transition\footnote{%
  A condition, bien s\^ur, de ne pas transformer la fonction de
  transition en lookup-table.}, %
mais l'augmentation du temps de calcul d\'epend de la taille de la
configuration initiale. %
Sur une machine SIMD l'augmentation de la taille du mot de voisinage
induit \'egalement une croissance lin\'eaire de la taille de l'espace et
du temps utilis\'e pour calculer une trajectoire cellulaire. %
Par contre l'augmentation de la taille d'espace d\'epend de la taille
du mot de voisinage et de la complexit\'e de la fonction de transition,
ce facteur de croissance de la taille d'espace est donc sup\'erieur \`a
celui d'une machine SISD\@. %
De plus la taille de l'architecture est fonction de la taille de la
configuration initiale, ce facteur pr\'esentant \'evidemment la source
principale des besoins de ressource spatiale. %
Par contre, pour prix de cette augmentation lin\'eaire des besoins
d'espace par rapport \`a la version SISD, le temps de calcul n'est plus
fonction de la taille de l'espace initial. %
Sur une machine cellulaire pure l'augmentation de la taille du mot de
voisinage induit une croissance exponentielle de la taille d'espace
via une croissance exponentielle de la taille de la table elle-m\^eme
facteur de la taille de la configuration initiale, le temps se
r\'eduisant \`a l'unit\'e. %

\section{Plan de lecture}

Dans le chapitre~\mygetref{space-time-law} nous pr\'esentons la
structure des objets de notre \'etude:~l'espace, le temps et les lois, %
ainsi que les m\'ethodes applicables sur ces objets:~\'edition, r\'eduction,
transformation, visualisation pour l'espace et le temps, plus
construction, alt\'eration, composition, m\'elange et s\'eparation pour les
lois. %
Nous pr\'esentons \'egalement dans ce chapitre l'architecture de nos
moteurs cellulaires. %

Dans le chapitre~\mygetref{rule-space} nous pr\'esentons les outils
conceptuels utilis\'es pour la quantification de l'espace des r\`egles, et
nous proposons un nouveau param\`etre de contr\^ole pour l'exploration et
la structuration de cet espace:~le param\`etre d'ench\^assement. %

Le chapitre~\mygetref{about-anneal} est consacr\'e \`a une \'etude
exp\'erimentale de la r\`egle~\arule{Anneal} bas\'ee sur l'exploitation de
nos outils. %
Les exp\'eriences r\'ealis\'ees nous am\`enent \`a proposer l'existence d'une
relative invariance spatio-temporelle des trajec\-toires associ\'ees \`a
une transition de phase dans l'espace des configurations initiales. %

La conclusion introduit les limites de notre approche qui \`a associ\'e
un constructivisme bottom-up \`a une \aquote{gestion de tables
  g\'en\'eralis\'ee}. %
Nous \'evoquons pour finir une approche alternative de r\'ealisation
cellulaire purement algorithmique. %

\mydumpfull{anneal-sum-1k-8k}%
{Une superposition de r\'eduction int\'egrale}%
{Une superposition de r\'eduction int\'egrale de \arule{Anneal}}%
{Cette image est obtenue par la superposition de la r\'eduction int\'egrale
  $\sum_{i=0}^{i=n}t_{i}xy$ des~$2^{10}$ et~$2^{13}$ premi\`eres
  g\'en\'erations d'un espace initial de taille~\asquare{1024} et de
  densit\'e\,\Sundemi.}

\mydumpfull{anneal-sum-1k-8k-xv}%
{}%
{Une autre colormap pour la figure
  \protect\mygetdump{anneal-sum-1k-8k}}%
{Une colormap en bandes. %
  Pour cette image, comme pour toutes les images en niveaux de gris
  que nous avons pr\'esent\'ees, un p\'eriph\'erique de sortie capable de
  restitution nuanc\'ee nous fait d\'efaut.}

%%
%% Local Variables:
%% mode: auto-fill
%% iso-accents-minor-mode: nil
%% TeX-command-default: "mytex"
%% TeX-master: "top"
%% Local IspellParsing: "latex-mode ~latin1"
%% Local IspellPersDict: "~/.ispell_lisp"
%% LocalWords: "l'auto-reproduction s\'equencement CA LAB cellsim SunView XView"
%% LocalWords: "Felicity scamper pattern cellang d'AC SIMD wafer scale l'IA MIMD"
%% LocalWords: "matching holographique multim\'edia multi-calculateur ad hoc WAN"
%% LocalWords: "Wide Area Network t\'el\'ecommunication bottom-up von Neumann cpu PE"
%% LocalWords: "adressable repartition-du-calcul l'overhead uds NEW DSP bitblt"
%% LocalWords: "morphog\'en\'etique g\'en\'ericit\'e cardinalit\'e fractales g\'en\'eticit\'e uvre"
%% LocalWords: "it\'er\'ee ergodique analytique-synth\'etique param\'etr\'e Neuman Turing"
%% LocalWords: "auto-reproducteur impl\'ement\'ee Conway univer GAPP it\`ere it\'erant"
%% LocalWords: "lookup-table mapping ind\'eterminisme indiff\'erentiable Langton DFA"
%% LocalWords: "\'epig\'en\'etique morphog\'en\`eses espace-temps Jaqueline Signorini past"
%% LocalWords: "auto-reproductrice monoprocesseur l'\'emulateur impl\'ementation"
%% LocalWords: "stricto sensu SISD CMOS approxim\'e space-time-law rule-space in"
%% LocalWords: "about-anneal Anneal constructivisme"
%% End:
